\chapter{Bipolar Medicines}

\section*{Definition and Overview}
Bipolar medicines are the medications used to treat bipolar disorder, a condition characterised by severe mood swings between mania and depression. The main groups are mood stabilisers, antipsychotics, and antidepressants. Most people with bipolar disorder require long-term treatment to prevent relapses, and the choice of medicine depends on the phase of the illness, previous responses, side effects, and individual health needs.

All of these medicines are prescribed by a doctor and require regular monitoring. When taken consistently and reviewed properly, they substantially reduce the frequency and severity of mood episodes and allow most people to live stable, productive lives.

\section*{Epidemiology}
Bipolar disorder affects roughly 1 to 2 per cent of the population, and almost all affected people are offered medication at some point.

\begin{itemize}
    \item \textbf{Sex distribution:} men and women are affected about equally
    \item \textbf{Age at treatment:} most people begin medication in their late teens or twenties, close to the age of onset of the disorder
    \item \textbf{Adherence:} many people stop taking a mood stabiliser within the first year, often because of side effects or a desire to feel normal, which is a major factor in relapse
    \item \textbf{Coexisting conditions:} substance use, anxiety, and physical illness are common and influence how medicines are chosen
\end{itemize}

\section*{Aetiology and Pathophysiology}
Bipolar medicines act on the brain systems believed to be disturbed in bipolar disorder.

\begin{itemize}
    \item \textbf{Lithium:} the oldest mood stabiliser; it dampens excessive signalling in brain pathways and reduces both manic and depressive relapse, although its exact mechanism remains incompletely understood
    \item \textbf{Anticonvulsants:} valproate and lamotrigine, originally developed for epilepsy, stabilise the electrical activity of nerve cells and help prevent mood swings
    \item \textbf{Antipsychotics:} such as olanzapine, quetiapine, and aripiprazole, block dopamine and serotonin receptors, calming overactive brain circuits during mania and psychosis
    \item \textbf{Antidepressants:} increase the availability of serotonin and noradrenaline to lift mood, but are used cautiously because they may trigger mania in some people
    \item \textbf{Benzodiazepines:} used briefly for severe agitation or insomnia, but not recommended as long-term treatment
\end{itemize}

\section*{Clinical Features}
Medicines are matched to the symptoms they are intended to target.

\begin{itemize}
    \item \textbf{During mania:} antipsychotics, mood stabilisers, and sometimes short-term benzodiazepines calm agitation, racing thoughts, and risky behaviour
    \item \textbf{During depression:} mood stabilisers, lamotrigine, quetiapine, or selected antidepressants lift low mood and restore energy and interest
    \item \textbf{Prevention:} lithium, valproate, lamotrigine, and some antipsychotics are taken continuously to reduce the frequency and severity of future episodes
    \item \textbf{Sleep and anxiety:} medicines may be used to restore normal sleep, which itself helps stabilise mood
    \item \textbf{Psychotic symptoms:} antipsychotics are used when delusions or hallucinations are present
\end{itemize}

\section*{Differential Diagnosis}
The decision to start a particular medicine depends on the correct diagnosis.

\begin{itemize}
    \item Unipolar depression, which does not usually require a mood stabiliser
    \item Cyclothymic disorder, where the treatment approach may differ
    \item Schizoaffective disorder and schizophrenia, which are treated mainly with antipsychotics
    \item Substance-induced mood disturbance, which may resolve when the substance is withdrawn
    \item Medical conditions such as thyroid disease, which require specific physical treatment rather than psychiatric medication
\end{itemize}

\section*{When to Seek Medical Care}
See a doctor promptly if a person with bipolar disorder has worsening mood symptoms, new side effects, or concerns about their medication. Never stop a mood stabiliser suddenly, because this can trigger a relapse; any planned change should be discussed with the prescribing doctor.

\textbf{Call 000 (or local emergency services) for:}
\begin{itemize}
    \item Thoughts of suicide or self-harm
    \item Severe mania with dangerous or reckless behaviour
    \item Confusion, seizures, or a serious adverse reaction to a medicine
    \item Inability to care for oneself, or severe agitation
\end{itemize}

\section*{Diagnosis and Investigations}
Before starting or adjusting medicines, doctors carry out baseline assessments and arrange ongoing monitoring.

\begin{itemize}
    \item \textbf{Mental state assessment:} to determine the current mood phase and guide the choice of medicine
    \item \textbf{Blood tests:} baseline kidney and thyroid function before lithium; liver function, full blood count, and a pregnancy test before valproate; and blood sugar and lipid levels before some antipsychotics
    \item \textbf{Electrocardiogram (ECG):} a heart rhythm check before some medicines, particularly in older people or those with heart disease
    \item \textbf{Lithium levels:} regular blood tests to keep the dose in a safe and effective range
    \item \textbf{Weight and metabolic monitoring:} regular checks of weight, waist circumference, blood sugar, and lipids during antipsychotic treatment
    \item \textbf{Pregnancy testing:} essential, because valproate and some other medicines can harm a developing baby
\end{itemize}

\section*{Management}
Treatment is individualised and reviewed regularly.

\begin{itemize}
    \item \textbf{Lithium:} effective for preventing both mania and depression; it requires dose adjustment and regular blood monitoring
    \item \textbf{Valproate:} useful for mania and rapid cycling, but generally avoided in women of childbearing age because of the risk of birth defects
    \item \textbf{Lamotrigine:} particularly helpful for preventing depressive episodes
    \item \textbf{Antipsychotics:} used for acute episodes and as maintenance treatment in many cases
    \item \textbf{Antidepressants:} reserved for depressive episodes and combined with a mood stabiliser in most cases
    \item \textbf{Combined treatment:} many people need more than one medicine to achieve stability
    \item \textbf{Psychological care alongside medicines:} psychoeducation and psychological therapy improve adherence and outcomes
\end{itemize}

\section*{Complications}
\begin{itemize}
    \item \textbf{Lithium side effects:} thirst, tremor, weight gain, underactive thyroid, and effects on kidney function with long-term use
    \item \textbf{Antipsychotic side effects:} weight gain, drowsiness, raised blood sugar, and movement problems
    \item \textbf{Anticonvulsant side effects:} sedation, dizziness, rash with lamotrigine, and effects on the liver with valproate
    \item \textbf{Antidepressant-induced mania or rapid cycling:} in susceptible people
    \item \textbf{Birth defects:} if valproate or lithium is taken during pregnancy
    \item \textbf{Drug interactions:} with other medicines, including some pain relievers, antibiotics, and heart medications
    \item \textbf{Withdrawal or relapse:} if medicines are stopped suddenly or without supervision
\end{itemize}

\section*{Prognosis and Outcomes}
When taken regularly, mood stabilisers substantially reduce the frequency and severity of mood episodes in many people with bipolar disorder, and lithium has been associated with a reduced risk of suicide in some studies.

However, side effects and the demands of monitoring mean that many people struggle to stay on treatment, and non-adherence is the most common cause of relapse. Outcomes are best when medicines are combined with psychological support, regular medical review, and a stable daily routine. Some people need changes of medicine over time as their illness and circumstances evolve.

\section*{Prevention and Self-Care}
\begin{itemize}
    \item Take medicines at the same time each day and use reminders or a dosette box
    \item Never stop or change a medicine without medical advice
    \item Keep all blood-test appointments for monitoring
    \item Report side effects early rather than stopping the medicine yourself
    \item Maintain regular sleep, exercise, and mealtimes
    \item Avoid alcohol and recreational drugs, which can destabilise mood and interact with medicines
    \item Carry a current list of medicines and tell every doctor and pharmacist what you take
    \item Discuss family planning openly with your doctor, because some medicines affect pregnancy
\end{itemize}

\section*{Sources and Further Reading}
The following organisations provide reliable information on medicines used in bipolar disorder for patients, families, and health professionals.

\begin{itemize}
    \item NHS. \textit{Information on bipolar disorder medicines and how they are used.} https://www.nhs.uk/
    \item NICE. \textit{Clinical guidance on the pharmacological management of bipolar disorder.} https://www.nice.org.uk/
    \item MedlinePlus. \textit{Information on mood stabilisers and related medicines.} https://medlineplus.gov/
    \item Mayo Clinic. \textit{Overview of bipolar disorder treatment, including medicines.} https://www.mayoclinic.org/
    \item National Institute of Mental Health. \textit{Educational resources on bipolar disorder treatments.} https://www.nimh.nih.gov/
\end{itemize}
